\documentclass[11pt]{article}
\usepackage{amsmath, amssymb, amsthm, graphicx, hyperref}

\title{Coherent Manifold Construction for Fictional Universes:\\
A Mathematical Framework Using Altitude Models, Spectral Geometry,\\
and High-Dimensional Conceptual Embeddings}

\author{Borealis S. Hedling}
\date{2026}

\begin{document}
\maketitle

\begin{abstract}
Fictional universes such as Dragon Ball, Star Wars, and Marvel exhibit 
complex internal structure: multiversal branching, nonlinear timelines, 
transformation hierarchies, cosmological domains, and narrative curvature. 
We propose a rigorous mathematical framework for modeling fictional 
universes as coherent high-dimensional manifolds. Our approach integrates 
the FiftyD Neutral Manifold, NDH Constellation symbolic geometry, the 
Hyperatlas narrative-construction protocol, and the A0--A10 altitude model 
for runtime physics. We show that fictional universes constitute rich, 
high-dimensional data sets suitable for spectral clustering, holonomy 
analysis, and conceptual embeddings. We introduce the repaired 8B 
Resonance Algebra (v2.0) as the spectral-temporal signal layer enabling 
coherent manifold construction. A Python toolkit is provided for 
constructing conceptual embeddings, saga submanifolds, holonomy curves, 
and cosmological fiber bundles. Dragon Ball is used as a case study.
\end{abstract}

\section{Introduction}

Fictional universes are often treated as narrative artifacts rather than 
mathematical objects. However, modern transmedia franchises exhibit 
structural complexity comparable to scientific domains: multiverses, 
retcons, nonlinear timelines, transformation operators, cosmological 
hierarchies, and cross-medium continuity constraints. These features 
suggest that fictional universes can be modeled as high-dimensional 
conceptual manifolds with well-defined geometry, topology, and operator 
algebra.

This preprint formalizes the mathematical foundations required to build 
coherent manifolds for fictional universes. We integrate four conceptual 
primitives:

\begin{enumerate}
  \item FiftyD Neutral Manifold: a 50-dimensional neutral coordinate space.
  \item Hyperatlas Protocol: narrative geometry, manifold expansion, and 
        trauma-informed stability.
  \item NDH Constellation: symbolic surfaces, nodes, and relational geometry.
  \item A0--A10 Altitude Model: geometric, spectral, modal, dual, and 
        non-dual runtime physics.
\end{enumerate}

We show how these primitives combine to produce coherent fictional 
manifolds and demonstrate their application to Dragon Ball.

\section{Mathematical Foundations}

\subsection{Neutral High-Dimensional Manifold}

Define the neutral manifold:
\[
M_{50} = \mathbb{R}^{50}
\]
with neutral metric:
\[
g_{ij} = \delta_{ij}.
\]

This manifold provides a bias-free embedding space for conceptual entities.

\subsection{NDH Constellation Geometry}

Define constellation:
\[
C = (N, S, E)
\]
where $N$ are nodes, $S$ are surfaces, and $E$ are edges.

Mapping:
\[
\iota : C \to M_{50}
\]
embeds symbolic geometry into conceptual space.

\subsection{Hyperatlas Narrative Geometry Field}

Define narrative geometry field:
\[
\Phi : M_{50} \to \mathbb{R}^k
\]
encoding adjacency, continuity, trauma-informed stability, and curvature.

\subsection{Altitude Model}

Altitude model:
\[
\mathcal{A} = \{A0, A1, \dots, A10\}
\]
provides operator semantics for conceptual physics.

\section{8B Resonance Algebra v2.0}

We define the repaired resonance field:
\[
R : M_{50} \to \mathbb{R}^3
\]
with:
\[
R(p) = (f(p), A(p), \phi(p)).
\]

We incorporate:
\begin{itemize}
  \item spectral continuity,
  \item stability envelopes,
  \item modal lift,
  \item kernel invariants,
  \item holonomy bounds,
  \item collapse surfaces,
  \item dual reprojection.
\end{itemize}

This algebra provides the spectral-temporal physics for fictional universes.

\section{Fictional Universes as High-Dimensional Data Sets}

We argue that fictional universes satisfy criteria for high-dimensional 
scientific data sets:

\begin{enumerate}
  \item \textbf{Multiversality}: multiple universes form fiber bundles.
  \item \textbf{Nonlinear timelines}: holonomy loops and retcon curvature.
  \item \textbf{Transformation operators}: spectral lift and modal adjacency.
  \item \textbf{Cosmological hierarchy}: topological stratification.
  \item \textbf{Narrative curvature}: trauma-informed stability fields.
\end{enumerate}

Dragon Ball is particularly rich due to:
\begin{itemize}
  \item transformation lineage (SSJ1 $\to$ SSJ Blue $\to$ UI),
  \item multiversal cosmology (12 universes),
  \item nonlinear timelines (Future Trunks),
  \item divine domains (Kai, Hakaishin),
  \item retcon curvature (Super vs Z).
\end{itemize}

\section{Python Toolkit}

We provide a Python toolkit for constructing fictional manifolds:

\begin{verbatim}
from fiftyd import NeutralManifold
from ndh import Constellation
from hyperatlas import NarrativeField
from altitude import AltitudeRuntime
from resonance import ResonanceField

M = NeutralManifold(dim=50)
C = Constellation.load("dbz_entities.json")
Phi = NarrativeField(M, C)
A = AltitudeRuntime()
R = ResonanceField(M)
\end{verbatim}

Toolkit supports:
\begin{itemize}
  \item conceptual embeddings,
  \item saga submanifolds,
  \item holonomy curves,
  \item cosmological fiber bundles,
  \item transformation operators.
\end{itemize}

\section{Case Study: Dragon Ball}

We embed:
\begin{itemize}
  \item characters,
  \item forms,
  \item sagas,
  \item timelines,
  \item universes,
  \item cosmology.
\end{itemize}

We construct:
\begin{itemize}
  \item transformation manifold,
  \item saga DAG,
  \item timeline holonomy,
  \item cosmological topology.
\end{itemize}

\section{Conclusion}

We provide a rigorous mathematical framework for fictional universe 
manifold construction. Dragon Ball demonstrates the viability of this 
approach. Future work includes A8 congruence geometry, A9 quantum-inspired 
runtime, and A10 conceptual operator physics.

\section{References}

\begin{itemize}
  \item Bostrom, N. (2014). \textit{Superintelligence}. Oxford University Press.
  \item Lakoff, G. (1987). \textit{Women, Fire, and Dangerous Things}. University of Chicago Press.
  \item Tegmark, M. (2014). \textit{Our Mathematical Universe}. Knopf.
  \item Zalta, E. (2021). \textit{Stanford Encyclopedia of Philosophy: Abstract Objects}.
\end{itemize}

\end{document}
